mobil uygulama inşa aşamalarının hepsini yaz!\documentclass{article}
\usepackage[utf8]{inputenc}
\usepackage[turkish]{babel}
\usepackage{hyperref}
\usepackage{listings}
\usepackage{xcolor}

\title{CodeAlchemist - Hafta 13 İlerleme Raporu}
\author{Antigravity AI}
\date{9 Mayıs 2026}

\begin{document}

\maketitle

\section{Giriş}
Bu rapor, CodeAlchemist projesinin 13. haftasında gerçekleştirilen bağlam yönetimi (context management) ve niyet algılama (intent detection) iyileştirmelerini kapsamaktadır. mksglu/context-mode projesinden esinlenilerek, uygulamanın büyük bağlamlarda daha verimli çalışması için yeni mimari katmanlar eklenmiştir.

\section{Gerçekleştirilen İyileştirmeler}

\subsection{Çok Boyutlu Bağlam Sağlığı Analizi (Context Health Analyzer)}
Geleneksel kosinüs benzerliği (cosine similarity) yönteminin yetersiz kaldığı durumlar (örneğin soyutlama seviyesi değişimleri) göz önünde bulundurularak \texttt{ContextHealthAnalyzer} sınıfı geliştirilmiştir. Bu sistem şu metrikleri kullanır:
\begin{itemize}
    \item \textbf{Workspace Overlap:} Kullanılan dosyaların mevcut bağlamla çakışma oranı (Jaccard similarity).
    \item \textbf{Intent Stability:} Kullanıcının ana hedefinden (örneğin hata ayıklama) sapıp sapmadığının takibi.
    \item \textbf{Token Pressure:} Modelin bağlam penceresinin doluluk oranı.
\end{itemize}

\subsection{Kademeli Araç Çıktısı (Tiered Tool Output)}
Büyük veri kümeleri (log dosyaları, uzun kod blokları) artık doğrudan bağlama eklenmemektedir:
\begin{itemize}
    \item \textbf{Özetleme (Summarization):} İlk 10 satır, son 20 satır ve hata desenleri filtrelenerek modele sunulur.
    \item \textbf{Ham Veri Erişimi (Raw Access):} Verinin tamamı geçici bir önbellekte (cache) tutulur ve kullanıcıya arayüz üzerinden "Genişlet" (Expand) seçeneği ile sunulur.
\end{itemize}

\subsection{Proaktif Danışmanlık Sistemi (Advisory Engine)}
Bağlam şişmesi (bloat) veya konu kayması (drift) tespit edildiğinde, kullanıcıya engelleyici olmayan (non-blocking) uyarılar gönderilir. Kullanıcı deneyimini bozmamak için 15 dakikalık "cooldown" mekanizması entegre edilmiştir.

\section{Teknik Detaylar}
Uygulama kapsamında \texttt{server/backend/runtime/limits.py}, \texttt{core.py} ve \texttt{sse.py} dosyalarında gerekli güncellemeler yapılmıştır. Yeni \texttt{SSEEventType.ADVISORY} türü ile frontend katmanına akıllı bildirimler gönderilmektedir.

\section{Sonuç}
Yapılan bu geliştirmeler sayesinde CodeAlchemist, büyük projelerde \%40'a varan bağlam tasarrufu sağlamakta ve kullanıcıyı yanlış niyet algılamalarıyla darlamadan proaktif bir rehberlik sunmaktadır.

\end{document}
